\chapter{Notazione Matematica e Glossario}

Per garantire il rigore formale e la chiarezza dell'esposizione, i capitoli successivi utilizzano una notazione matematica derivata dalla teoria dei tipi e dalla semantica funzionale. Questo capitolo funge da riferimento rapido per la simbologia e i concetti formali introdotti nella formalizzazione del sistema.

\section{Convenzioni Notazionali}
Il modello utilizza le seguenti convenzioni standard per descrivere le strutture dati e le funzioni del sistema:

\begin{itemize}
    \item \textbf{Tuple (Prodotto Cartesiano)} $\langle x_1, x_2, \dots, x_n \rangle$: Rappresentano tipi composti (corrispondenti alle \texttt{struct} o classi nel codice sorgente). Ogni elemento della tupla definisce un campo o una proprietà essenziale dell'entità.
    \item \textbf{Unione Disgiunta (Tipo Somma)} $A \mid B \mid C$: Indica che un valore può appartenere a uno e un solo tipo tra quelli elencati (corrispondente agli \texttt{enum} nei linguaggi funzionali o in Rust). È utilizzato per modellare il polimorfismo o l'incertezza.
    \item \textbf{Vettori e Sequenze} $\vec{X}$: Indica una sequenza finita e ordinata di elementi di tipo $X$ (corrispondente alle liste, array o \texttt{Vec} in memoria).
    \item \textbf{Opzionalità} $X^?$: Indica un valore che può essere presente (di tipo $X$) oppure assente (corrispondente a \texttt{Option<T>} o al valore \texttt{null}).
    \item \textbf{Firme Funzionali} $f : A \to B$: Denota una funzione $f$ che accetta un input dal dominio $A$ e restituisce un output nel codominio $B$. Se l'input è composto da più argomenti, si utilizza il prodotto cartesiano (es. $A \times B \to C$).
\end{itemize}

\section{Glossario dei Tipi e dei Domini}
Di seguito sono definiti i simboli principali che costituiscono l'alfabeto della nostra Rappresentazione Intermedia (IR) e della pipeline di analisi.

\subsection{Tipi Base e Identificatori}
\begin{itemize}
    \item $\mathcal{I}$ (\textbf{Identifier}): L'insieme di tutte le stringhe alfanumeriche valide utilizzate per nominare variabili, funzioni o classi.
    \item $\mathcal{QN}$ (\textbf{Qualified Name}): Una sequenza gerarchica di identificatori ($\langle i_1, \dots, i_n \rangle$) che rappresenta il percorso assoluto e univoco di un'entità (es. il namespace o il modulo).
    \item $\tau$ (\textbf{Type Reference}): Un riferimento astratto a un tipo. Può rappresentare un tipo primitivo, un tipo definito dall'utente (attraverso un $\mathcal{QN}$), o un tipo temporaneamente "sconosciuto" (\texttt{Unknown}) in attesa di risoluzione.
\end{itemize}

\subsection{Membri Strutturali}
\begin{itemize}
    \item $F$ (\textbf{Field}): Un campo dati appartenente a una struttura, definito dalla coppia identificatore e tipo ($\langle \mathcal{I}, \tau \rangle$).
    \item $P$ (\textbf{Parameter}): Un parametro in ingresso a una funzione, comprendente nome opzionale, tipo e indicatore di variadicità.
    \item $M$ (\textbf{Method}): Una funzione legata allo stato di un tipo strutturato (metodo di classe/struct).
    \item $FN$ (\textbf{Free Function}): Una funzione globale definita a livello di modulo, indipendente da qualsiasi tipo strutturato.
\end{itemize}

\subsection{Entità Architetturali}
\begin{itemize}
    \item $\mathcal{S}$ (\textbf{Structured Type}): L'astrazione di classi, struct, interfacce e trait. Aggrega campi, metodi, gerarchie di ereditarietà e tipi annidati.
    \item $\mathcal{B}_{lk}$ (\textbf{Implementation Block}): Costrutto tipico di linguaggi come Rust o C++ per separare la dichiarazione dei dati dalla definizione dei comportamenti. Associa un insieme di metodi a un tipo target.
    \item $\mathcal{C}$ (\textbf{Component}): Il tipo somma universale che abbraccia qualsiasi entità analizzabile dal sistema (Moduli, Tipi Strutturati, Funzioni Libere, Blocchi di Implementazione).
    \item $\mathcal{M}$ (\textbf{Module}): L'unità gerarchica fondamentale. Il modulo radice funge da dominio della Rappresentazione Intermedia ($\mathcal{D}_{IR}$) e aggrega ricorsivamente tutto il codice sorgente analizzato.
\end{itemize}

\subsection{Modello di Risoluzione e Grafo}
\begin{itemize}
    \item $\Sigma$ (\textbf{Symbol Table}): La funzione di mappa globale che associa i nomi qualificati ($\mathcal{QN}$) alle effettive istanze dei componenti ($\mathcal{C}$).
    \item $\Gamma$ (\textbf{Resolution Context}): Il contesto locale (comprendente scope attuale e importazioni) utilizzato per risolvere riferimenti ambigui durante la fase di risoluzione.
    \item $\mathcal{G}$ (\textbf{Dependency Graph}): Il grafo orientato finale $\langle V, E \rangle$. I vertici $V$ sono i componenti estratti dal codice e gli archi $E$ sono le relazioni semantiche di dipendenza (ereditarietà, composizione, uso).
\end{itemize}

\subsection{Fasi dell'Analisi}
\begin{itemize}
    \item $\Phi$ (\textbf{Full Analysis}): La funzione matematica complessiva che rappresenta l'intero programma, calcolata come la composizione delle tre fasi sottostanti ($\Phi = G \circ R \circ E$).
    \item $E$ (\textbf{Extraction}): La prima fase dipendente dal linguaggio che trasforma i nodi dell'albero sintattico nel Modulo Radice $\mathcal{M}$.
    \item $R$ (\textbf{Resolution}): La fase agnostica che naviga la gerarchia $\mathcal{M}$ risolvendo e validando tutti i riferimenti $\tau$ grazie al contesto $\Gamma$.
    \item $G$ (\textbf{Graphing}): L'ultima fase che converte il modello a oggetti risolto nel grafo architetturale $\mathcal{G}$.
\end{itemize}
